% !TeX program = xelatex
% ============================================================
% research/report/preamble.tex — преамбула Отчёта исследовательского
% курсового проекта ОП «Программная инженерия» ФКН НИУ ВШЭ.
%
% Оформление по ГОСТ 7.32-2017 «Отчёт о НИР» (класс extreport — главы):
% поля 30/15/20/20, TNR 12 пт, интервал 1,5, абзац 1,25 см, сквозная
% нумерация рисунков/таблиц. ЧИСТЫЙ документ — БЕЗ штампа/рамки/листа
% утверждения (в отличие от ЕСПД-документов курсового). Язык интерфейса
% разделов — по project.lang (RU/EN). Данные — из common/meta (\hse…).
% ============================================================
\newcommand{\hseDefaultMono}{Consolas}
\input{../common/fonts}

\usepackage[russian,english]{babel}


% Поля страницы — по настройке проекта report_margins (Настройки → Поля
% страницы Отчёта). По умолчанию ГОСТ 7.32-2017 (30/15/20/20); форма Прил. 10
% методички ВШЭ — 25/10/20/20. Переключение не трогает исходники.
((* if 'Прил' in (project.meta.report_margins | default('')) *))
\usepackage[a4paper,left=25mm,right=10mm,top=20mm,bottom=20mm]{geometry}
((* else *))
\usepackage[a4paper,left=30mm,right=15mm,top=20mm,bottom=20mm]{geometry}
((* endif *))
\usepackage{setspace}
\onehalfspacing
\usepackage{indentfirst}
\setlength{\parindent}{1.25cm}
\setlength{\parskip}{6pt}

\usepackage{microtype}
\usepackage{fvextra}
\usepackage{csquotes}
\usepackage{amsmath}
\usepackage{amssymb}
\usepackage{graphicx}
\usepackage{letltxmacro}
\LetLtxMacro{\oldincludegraphics}{\includegraphics}
\renewcommand{\includegraphics}[2][]{%
  \IfFileExists{#2}{%
    \oldincludegraphics[#1]{#2}%
  }{%
    \IfFileExists{#2.png}{\oldincludegraphics[#1]{#2.png}}{%
      \IfFileExists{#2.jpg}{\oldincludegraphics[#1]{#2.jpg}}{%
        \fcolorbox{red}{white}{%
          \begin{minipage}{0.5\textwidth}\centering
            \textcolor{red}{\textbf{Файл не найден:}} \\ \texttt{\detokenize{#2}}
          \end{minipage}%
        }%
      }%
    }%
  }%
}
\usepackage{xcolor}
\usepackage{hyperref}
\usepackage{lastpage}
\usepackage{float}
\usepackage{array}
\usepackage{booktabs}
\usepackage{longtable}
\usepackage{tabularx}
\usepackage{adjustbox}
\usepackage{xurl}
\urlstyle{same}
\usepackage{hyphenat}
\usepackage{caption}
% Ссылки [n] печатаются в строку (ГОСТ 7.0.5-2008), не надстрочно.
\usepackage{cite}
\DeclareCaptionLabelSeparator{emdash}{~---~}
\captionsetup[table]{position=top,labelsep=emdash,justification=raggedright,singlelinecheck=false}
\captionsetup[figure]{position=bottom,labelsep=emdash}

% ── Листинги кода ───────────────────────────────────────────
\usepackage{listings}
\renewcommand{\lstlistingname}{Листинг}
\definecolor{codegray}{rgb}{0.5,0.5,0.5}
\definecolor{backcolour}{rgb}{0.95,0.95,0.92}
\lstset{
  basicstyle=\footnotesize\ttfamily,
  backgroundcolor=\color{backcolour},
  commentstyle=\color{codegray}\itshape,
  numberstyle=\tiny\color{codegray},
  breaklines=true, captionpos=b, numbers=left, numbersep=5pt,
  keepspaces=true, showspaces=false, showstringspaces=false, showtabs=false,
  tabsize=2, frame=single, rulecolor=\color{gray!30},
  xleftmargin=10pt, xrightmargin=10pt,
}

\usepackage{enumitem}
\setlist{nosep}
\setlist[itemize,1]{label=--}
\usepackage{chngcntr}
\counterwithout{table}{chapter}
\counterwithout{figure}{chapter}
\renewcommand{\thetable}{\arabic{table}}
\renewcommand{\thefigure}{\arabic{figure}}
\usepackage{fancyhdr}
\pagestyle{plain}

\usepackage{titlesec}
\usepackage{tikz}
\usetikzlibrary{shapes.geometric, arrows, positioning, calc, er, fit}

\emergencystretch=5em
\tolerance=2000
\hyphenpenalty=300
\exhyphenpenalty=300
\titleformat{\chapter}{\bfseries\centering}{\thechapter}{1em}{}
\titleformat{\section}{\bfseries}{\thesection}{1em}{}
\titlespacing*{\section}{1.25cm}{*2}{6pt}
\titleformat{\subsection}{\bfseries}{\thesubsection}{1em}{}
\titlespacing*{\subsection}{1.25cm}{*2}{6pt}
\titleformat{\subsubsection}{\bfseries}{\thesubsubsection}{1em}{}


\renewcommand{\contentsname}{Contents}

\setcounter{tocdepth}{3}
\setcounter{secnumdepth}{3}

\hypersetup{unicode=true,breaklinks=true,bookmarksopen=true,bookmarksnumbered=true,colorlinks=true,linkcolor=black,urlcolor=blue,citecolor=blue}

% Значки для матриц сравнения (обзор аналогов/подходов).
\newcommand{\cmpplus}{\textbf{+}}
\newcommand{\cmpminus}{\textbf{--}}
\newcommand{\cmpplanned}{\textcolor{blue}{$\pm$}}

% Формат номеров в списке источников с точкой.
\makeatletter
\renewcommand\@biblabel[1]{#1.}
\makeatother

% ── Данные проекта (\hse…), затем «человеческие» имена (commands) ──
% meta даёт \hseFill и все \hse-макросы; commands — \signdateblank и пр.,
% которые использует титул (common/commands НЕ тянет штамп — он в styles).
\input{../common/meta}
\input{../common/commands}
